\documentclass[11pt]{article}

\usepackage[letterpaper,margin=1in]{geometry}
\usepackage[T1]{fontenc}
\usepackage[utf8]{inputenc}
\usepackage{lmodern}
\usepackage{microtype}
\usepackage{amsmath,amssymb,mathtools,bm}
\usepackage{enumitem}
\usepackage{xcolor}
\usepackage{booktabs}
\usepackage{array}
\usepackage{hyperref}
\usepackage[most]{tcolorbox}
\usepackage{fancyhdr}
\usepackage{titlesec}
\usepackage{parskip}
\usepackage{needspace}

\definecolor{qftblue}{HTML}{245A8D}
\definecolor{qftgreen}{HTML}{2F6B4F}
\definecolor{qftred}{HTML}{A23B3B}
\definecolor{qftgray}{HTML}{F4F6F8}

\hypersetup{
  colorlinks=true,
  linkcolor=qftblue,
  urlcolor=qftblue,
  citecolor=qftgreen,
  pdftitle={PHYS 662 Lecture 7: Symmetries and Noether's First Theorem},
  pdfauthor={Nima Lashkari}
}

\pagestyle{fancy}
\fancyhf{}
\lhead{PHYS 662: Quantum Field Theory I}
\rhead{Lecture 7}
\cfoot{\thepage}
\setlength{\headheight}{14pt}

\titleformat{\section}{\Large\bfseries\color{qftblue}}{\thesection.}{0.6em}{}
\titleformat{\subsection}{\large\bfseries\color{qftgreen}}{\thesubsection}{0.6em}{}
\titleformat{\subsubsection}{\normalsize\bfseries\color{qftred}}{\thesubsubsection}{0.6em}{}

\setlist[itemize]{leftmargin=2em,itemsep=0.25em,topsep=0.4em}
\setlist[enumerate]{leftmargin=2em,itemsep=0.25em,topsep=0.4em}

\newcommand{\dd}{\mathrm{d}}
\newcommand{\R}{\mathbb{R}}
\newcommand{\Lag}{\mathcal{L}}
\newcommand{\cM}{\mathcal{M}}
\newcommand{\SO}{\mathrm{SO}}
\newcommand{\OO}{\mathrm{O}}

\newtcolorbox{keyidea}[1][]{
  enhanced, breakable, colback=qftgray, colframe=qftblue,
  boxrule=0.8pt, arc=1.5mm, title={#1}, fonttitle=\bfseries,
  coltitle=white
}
\newtcolorbox{remarkbox}[1][]{
  enhanced, breakable, colback=white, colframe=qftgreen,
  boxrule=0.7pt, arc=1.5mm, title={#1}, fonttitle=\bfseries,
  coltitle=white
}
\newtcolorbox{warningbox}[1][]{
  enhanced, breakable, colback=white, colframe=qftred,
  boxrule=0.7pt, arc=1.5mm, title={#1}, fonttitle=\bfseries,
  coltitle=white
}


\newcommand{\Th}{\Theta}
\newcommand{\Eul}{\mathcal E}
\newcommand{\Lie}{\mathscr L}
\begin{document}
\pagenumbering{gobble}
\begin{titlepage}
\centering
\vspace*{1.2in}
{\scshape\Large PHYS 662 --- Quantum Field Theory I\par}
\vspace{0.35in}
{\Huge\bfseries\color{qftblue} Lecture 7\par}
\vspace{0.15in}
{\LARGE\bfseries Symmetries and Noether's First Theorem\par}
\vspace{0.12in}
{\Large\color{qftgreen} From Transformations to Conserved Currents\par}
\vspace{0.6in}
{\large Nima Lashkari\par}
{\large Fall 2026\par}
\vfill
\begin{minipage}{0.86\textwidth}
\small
A continuous symmetry of a local action produces a local conservation law.
We establish this statement by carefully comparing fields at the same
coordinate and at displaced points, varying both the integrand and the
integration region, and retaining the boundary term. Translations give
energy and momentum, internal phase rotations give a charge, and Lorentz
transformations give angular momentum and boost charges. The examples also
explain the roles of boundary flux, current conventions, and field spin.

\medskip
Course page: \url{https://nimalashkari.com/quantum-field-theory-2026-part-i/}
\end{minipage}
\vspace{0.5in}
\end{titlepage}
\pagenumbering{roman}
\tableofcontents
\newpage
\pagenumbering{arabic}

\section*{Roadmap}
\addcontentsline{toc}{section}{Roadmap}
Lecture 5 derived the field equations from stationary action. Lecture 6
provided actions on which to practice. We now ask what additional
information follows when an action has a continuous symmetry. The answer
is a current whose divergence vanishes on every classical solution.

% We follow the sequence of the slides: spacetime maps and field variations;
% the first variation on a finite region; Noether's first theorem; and the
% examples of translations, global \(U(1)\), and Lorentz transformations.
% The recurring distinction is between a symmetry identity, valid for
% arbitrary fields, and a conservation equation, valid on solutions.

\begin{keyidea}[Conventions]
We use \(D=d+1\), \(x^0=t\), \(\hbar=c=1\), and
\(\eta_{\mu\nu}=\operatorname{diag}(1,-1,\ldots,-1)\).
Greek indices label spacetime components; \(a,b\) label field components;
\(I\) labels independent symmetry parameters. Repeated indices are summed.
The action calculation uses Cartesian Minkowski coordinates.
All variation formulas retain only first order in the small parameters.
\end{keyidea}

\section{Spacetime maps, coordinates, and fields}

\subsection{A coordinate chart is a labeling of points}
A spacetime manifold \(\cM\) is a collection of events with a smooth local
coordinate structure. A chart is an invertible map
\begin{equation}
\alpha:U\subset\cM\longrightarrow\alpha(U)\subset\R^D,
\qquad p\longmapsto x=\alpha(p).
\end{equation}
A second chart \(\alpha'\) gives \(x'=\alpha'(p)\). On their common domain,
\begin{equation}
\boxed{x'=\alpha'\circ\alpha^{-1}(x).}
\label{eq:charts}
\end{equation}
The order matters: first recover the event from its old coordinates, then
assign its new coordinates. A chart is generally local; one need not be
able to label the entire manifold with a single chart.

A diffeomorphism \(f:\cM\to\cM\) is an invertible smooth map with a smooth
inverse. Unlike a change of chart, it can map one event \(p\) to a different
event \(f(p)\). In a fixed chart its coordinate expression is
\(F=\alpha\circ f\circ\alpha^{-1}\).

\subsection{Active and passive descriptions}
For an \emph{active} transformation, keep the chart fixed and move the
field configuration by \(f\). A scalar is transported with its value:
\begin{equation}
\widetilde\phi'(f(p))=\widetilde\phi(p).
\end{equation}
In coordinates, with \(x'=F(x)\), this reads
\begin{equation}
\boxed{\phi'(x')=\phi(x),\qquad
\phi'(x)=\phi(F^{-1}(x)).}
\label{eq:activeScalar}
\end{equation}
For \(F(x)=x+a\), a profile moves by \(+a\), and its new value at a fixed
coordinate is \(\phi'(x)=\phi(x-a)\). For example, the peak of
\(\exp[-(x-b)^2/(2\sigma^2)]\) moves from \(b\) to \(b+a\).
This provides a direct check of the inverse and the sign.

For a \emph{passive} transformation, keep the event and geometric field
fixed and change the chart. If \(x'=F(x)\), the two coordinate descriptions
again obey \(\phi'(x')=\phi(x)\). The same algebraic relation can therefore
appear in both descriptions, but the meanings of the primes differ.
A passive relabeling does not move the physical profile. Moreover, a passive
change to a frame displaced by \(+a\) commonly uses \(x'=x-a\), the inverse
coordinate map to an active displacement by \(+a\).

\begin{warningbox}[Evaluating a field elsewhere is not its active transformation]
The Taylor expansion \(\phi(x+\xi)=\phi(x)+\xi^\mu\partial_\mu\phi(x)\)
compares the original field at two arguments. With our active convention
\(x'=x+\xi\), the transformed scalar instead satisfies
\(\phi'(x)=\phi(x)-\xi^\mu\partial_\mu\phi(x)\) to first order.
One must specify which field and which argument are being compared.
\end{warningbox}

\subsection{Coordinate covariance and physical symmetry}
An isometry is a diffeomorphism preserving the metric:
\(f^*g=g\). In Minkowski space,
\begin{equation}
x'^\mu=\Lambda^\mu{}_{\nu}x^\nu+a^\mu,
\qquad \Lambda^T\eta\Lambda=\eta,
\end{equation}
is a Poincare transformation; the case \(a=0\) is a Lorentz transformation.
Infinitesimally an isometry has displacement \(\xi\) obeying
\begin{equation}
\partial_\mu\xi_\nu+\partial_\nu\xi_\mu=0.
\label{eq:Killing}
\end{equation}
This is the Killing equation in Cartesian flat coordinates.

Coordinate independence is a requirement on how we describe a theory.
It does not make every active diffeomorphism a symmetry of matter on a
fixed metric. A generic active deformation changes lengths and gradients
relative to that metric. Even an isometry can fail to be a symmetry if
an external source or boundary condition breaks it. In general coordinates
one writes \(\sqrt{|g|}\,\dd^D x\) and covariant derivatives; here we keep
\(\eta\) fixed and test the actual action and its backgrounds.

\section{Two field variations and two kinds of symmetry}

Consider an infinitesimal transformation of coordinates and fields,
\begin{equation}
x'^\mu=x^\mu+\xi^\mu(x),\qquad
\phi'^a(x')=\phi^a(x)+\Delta\phi^a(x).
\end{equation}
Define the variation at the same coordinate by
\begin{equation}
\delta\phi^a(x):=\phi'^a(x)-\phi^a(x).
\end{equation}
Taylor expansion of \(\phi'^a(x')\) gives
\begin{equation}
\boxed{\Delta\phi^a=\delta\phi^a+\xi^\mu\partial_\mu\phi^a,
\qquad \delta\phi^a=\Delta\phi^a-\xi^\mu\partial_\mu\phi^a.}
\label{eq:twoVariations}
\end{equation}
The derivative of \(\delta\phi\) multiplied by \(\xi\) would be second
order and has been dropped.

\begin{center}
\small
\begin{tabular}{@{}p{0.16\textwidth}p{0.34\textwidth}p{0.41\textwidth}@{}}
\toprule
Quantity & Comparison & Scalar transported by \(x'=x+\xi\) \\
\midrule
\(\Delta\phi\) & New field at \(x'\), old field at \(x\) & \(0\) \\
\(\delta\phi\) & New and old fields at the same argument &
\(-\xi^\mu\partial_\mu\phi\) \\
\bottomrule
\end{tabular}
\end{center}

An internal transformation leaves the spacetime argument fixed:
\(\xi^\mu=0\). For two real fields, an internal rotation is
\begin{equation}
\begin{pmatrix}\phi'^1\\\phi'^2\end{pmatrix}
=\begin{pmatrix}\cos\alpha&-\sin\alpha\\\sin\alpha&\cos\alpha\end{pmatrix}
\begin{pmatrix}\phi^1\\\phi^2\end{pmatrix}.
\end{equation}
A spacetime transformation has nonzero \(\xi\), and can also rotate tensor
or spinor components. Consequently, \(\delta\phi\) is \emph{not} synonymous
with an internal transformation: it contains the displacement contribution
as well. The component change \(\Delta\phi\) vanishes for a transported
scalar, but generally not for a vector or spinor.

For finitely many constant parameters \(\varepsilon^I\), write
\begin{equation}
\xi^\mu=\varepsilon^I X_I^\mu,
\qquad \Delta\phi^a=\varepsilon^I\Phi_I^a,
\qquad \delta\phi^a=\varepsilon^I\chi_I^a,
\quad \chi_I^a:=\Phi_I^a-X_I^\nu\partial_\nu\phi^a.
\label{eq:generators}
\end{equation}
The functions \(X_I,\Phi_I\) describe the infinitesimal transformation.
Here the coordinate map is a point transformation, with \(X_I^\mu(x)\).
The \(\varepsilon^I\) are independent constants, not arbitrary functions
of position. This is the setting of Noether's first theorem.

\section{Varying the action on a finite region}
\label{sec:variation}

\subsection{The fixed-coordinate first variation}
For a collection of fields take one local Lagrangian,
\begin{equation}
S_B[\phi]=\int_B\dd^D x\,\Lag(\phi^a,\partial_\mu\phi^a,x).
\end{equation}
It may couple different components; it need not be a sum of independent
single-field Lagrangians. Complex scalars can be treated as pairs of real
fields, or varied using \(\phi\) and \(\phi^*\) as independent variables
before imposing complex conjugacy.

Define
\begin{equation}
\pi_a^\mu:=\frac{\partial\Lag}{\partial(\partial_\mu\phi^a)},
\qquad
\Eul_a:=\frac{\partial\Lag}{\partial\phi^a}-\partial_\mu\pi_a^\mu.
\label{eq:ELdefinitions}
\end{equation}
The equations of motion are \(\Eul_a=0\); this condition is called
\emph{on shell}. An identity valid without imposing it is \emph{off shell}.
Since the fixed-coordinate variation commutes with differentiation,
\begin{align}
\delta\Lag
&=\frac{\partial\Lag}{\partial\phi^a}\delta\phi^a
 +\pi_a^\mu\partial_\mu\delta\phi^a\notag\\
&=\Eul_a\delta\phi^a+\partial_\mu(\pi_a^\mu\delta\phi^a).
\label{eq:basicVariation}
\end{align}
In these expressions, \(\partial_\mu\) differentiates the full spacetime
function, including the fields' dependence on \(x\). It is the total
spacetime derivative denoted \(D_\mu\) on the slides. By contrast,
\((\partial\Lag/\partial x^\mu)_{\rm explicit}\) holds the field and
its derivatives fixed.

\subsection{The derivative and the volume element also transform}
At displaced arguments, the derivative transformation is
\begin{equation}
\Delta(\partial_\mu\phi^a)
=\partial_\mu\Delta\phi^a
 -(\partial_\mu\xi^\nu)\partial_\nu\phi^a.
\label{eq:derivativeVariation}
\end{equation}
This follows from
\(\partial'_\mu=(\delta_\mu{}^\nu-\partial_\mu\xi^\nu)\partial_\nu\).
Using Eq.~\eqref{eq:twoVariations}, it also equals
\(\partial_\mu\delta\phi^a+\xi^\nu\partial_\nu\partial_\mu\phi^a\).
Thus \(\Delta\) does not commute with a coordinate derivative in the same
way that \(\delta\) does.

Let \(B'=F(B)\), and compare \(S_{B'}[\phi']\) to \(S_B[\phi]\).
Near the identity the Jacobian is positive, and
\begin{align}
\dd^D x'&=\det\left(\delta^\mu{}_{\nu}+\partial_\nu\xi^\mu\right)
\dd^D x\notag\\
&=\left(1+\partial_\mu\xi^\mu\right)\dd^D x+O(\varepsilon^2).
\label{eq:Jacobian}
\end{align}
For a matrix \(A\), the linear term of \(\det(I+\varepsilon A)\) is
\(\varepsilon\operatorname{tr}A\). This proves the Jacobian formula.
The finite-change definition uses an absolute determinant; an
infinitesimal map connected to the identity lets us omit the absolute sign.

Define \(\Delta\Lag\) by comparing the transformed integrand at \(x'\)
with the old integrand at \(x\). The chain rule gives
\begin{equation}
\Delta\Lag=\delta\Lag+\xi^\mu\partial_\mu\Lag.
\end{equation}
Pulling the transformed integral back to \(B\), its first variation is
\begin{align}
\Delta S_B
&=\int_B\dd^D x\left(\Delta\Lag+\Lag\partial_\mu\xi^\mu\right)\notag\\
&=\int_B\dd^D x\left[\delta\Lag+
\partial_\mu(\Lag\xi^\mu)\right].
\label{eq:movingRegion}
\end{align}
The term \(\Lag\partial_\mu\xi^\mu\) from the Jacobian is essential:
\((\partial_\mu\Lag)\xi^\mu\) alone is not the divergence of
\(\Lag\xi^\mu\).

\subsection{The boundary term contains the conserved current}
Combining Eqs.~\eqref{eq:basicVariation} and \eqref{eq:movingRegion},
\begin{equation}
\boxed{\Delta S_B=\int_B\dd^D x\,\Eul_a\delta\phi^a
+\int_{\partial B}\dd\Sigma_\mu
\left(\pi_a^\mu\delta\phi^a+\Lag\xi^\mu\right).}
\label{eq:master}
\end{equation}
Here \(\dd\Sigma_\mu\) is the oriented boundary element, defined so that
\(\int_B\partial_\mu V^\mu\,\dd^D x=
\int_{\partial B}V^\mu\dd\Sigma_\mu\).
On shell the bulk term vanishes, but the boundary term need not.

For the action principle, we used arbitrary variations vanishing at the
boundary to derive \(\Eul_a=0\). A symmetry transformation is a particular
variation that generally does not vanish on \(\partial B\). Its boundary
term must be retained. The key issue is the boundary condition on the
variation, not whether the boundary happens to be at infinity.

Introduce the canonical energy-momentum tensor
\begin{equation}
\boxed{\Th^\mu{}_{\nu}
=\pi_a^\mu\partial_\nu\phi^a-\delta^\mu{}_{\nu}\Lag.}
\label{eq:Theta}
\end{equation}
Then the boundary integrand becomes
\begin{equation}
\pi_a^\mu\delta\phi^a+\Lag\xi^\mu
=\pi_a^\mu\Delta\phi^a-\Th^\mu{}_{\nu}\xi^\nu.
\label{eq:boundaryTheta}
\end{equation}
The index \(\mu\) labels the direction of current flow; \(\nu\) labels
which translation charge is transported. The displacement multiplies
\(\Th\) in Eq.~\eqref{eq:boundaryTheta}; it is not part of the definition
of \(\Th\) itself.

\section{Noether's first theorem}

\subsection{A local symmetry condition}
A variational symmetry is a transformation for which the action density,
including its measure, changes by at most a total derivative. For each
constant parameter we require the off-shell identity
\begin{equation}
\boxed{\Delta\Lag+\Lag\partial_\mu\xi^\mu
=\varepsilon^I\partial_\mu K_I^\mu.}
\label{eq:symmetryCondition}
\end{equation}
The local functions \(K_I^\mu\) describe an allowed boundary contribution.
Strict invariance corresponds to \(K_I^\mu=0\).
For a finite region this means
\(\Delta S_B=\varepsilon^I\int_{\partial B}K_I^\mu\dd\Sigma_\mu\).
The full action is invariant when the relevant boundary term vanishes,
or the boundary formulation is chosen to accommodate it.

\begin{keyidea}[Noether's first theorem]
Each independent continuous global variational symmetry of a local action
has a local current whose divergence vanishes on solutions of the
Euler--Lagrange equations. The symmetry condition is an off-shell identity;
current conservation is its on-shell consequence.
\end{keyidea}
A discrete transformation has no infinitesimal parameter and is outside
this statement. Gauge transformations with arbitrary spacetime-dependent
parameters involve Noether's second theorem and identities among the
field equations; they do not supply an independent ordinary bulk charge
for every function used as a gauge parameter. A Noether current can also
be trivial or differ from another representative by a boundary term.

\subsection{Proof and the current formula}
Substituting Eq.~\eqref{eq:generators} into the local form of
Eq.~\eqref{eq:master}, and comparing with Eq.~\eqref{eq:symmetryCondition},
gives, for each \(I\),
\begin{equation}
\Eul_a\chi_I^a+
\partial_\mu\left(\pi_a^\mu\chi_I^a+\Lag X_I^\mu-K_I^\mu\right)=0.
\end{equation}
Thus define
\begin{equation}
\boxed{J_I^\mu
=\pi_a^\mu\chi_I^a+\Lag X_I^\mu-K_I^\mu
=\pi_a^\mu\Phi_I^a-\Th^\mu{}_{\nu}X_I^\nu-K_I^\mu.}
\label{eq:NoetherCurrent}
\end{equation}
The complete identity is
\begin{equation}
\boxed{\partial_\mu J_I^\mu=-\Eul_a\chi_I^a
\quad\Longrightarrow\quad\partial_\mu J_I^\mu=0\ \text{on shell}.}
\label{eq:NoetherIdentity}
\end{equation}
Equivalently, on a solution the flux of \(J_I\) through the boundary of
every sufficiently small region vanishes. Since the region is arbitrary,
its local divergence must vanish. Checking only that one integral over
all spacetime is zero would not establish this local conclusion.

\subsection{An equivalent calculation at fixed coordinates}
If we choose to keep the integration region fixed, use only
\(\delta\phi^a=\varepsilon^I\chi_I^a\). The symmetry condition becomes
\begin{equation}
\delta\Lag=\varepsilon^I\partial_\mu F_I^\mu,
\qquad F_I^\mu=K_I^\mu-\Lag X_I^\mu.
\end{equation}
The familiar formula is then
\begin{equation}
J_I^\mu=\pi_a^\mu\chi_I^a-F_I^\mu.
\end{equation}
It agrees with Eq.~\eqref{eq:NoetherCurrent}. In particular, a spacetime
symmetry can be strictly invariant when the region is moved, while its
fixed-coordinate Lagrangian changes by a total derivative. These are two
consistent ways to do the same calculation.

As a useful diagnostic, temporarily replace a constant parameter by a
smooth function \(\varepsilon(x)\) in the fixed-coordinate variation.
For the point transformations considered here, Eq.~\eqref{eq:basicVariation}
and the constant-parameter symmetry identity imply
\begin{equation}
\delta_{\varepsilon(x)}\Lag
=\partial_\mu(\varepsilon F^\mu)
+(\partial_\mu\varepsilon)J^\mu.
\label{eq:localParameter}
\end{equation}
Thus the coefficient of \(\partial_\mu\varepsilon\), after isolating a
boundary term, is the Noether current. For compactly supported
\(\varepsilon(x)\), stationarity on shell yields
\(-\int\varepsilon\partial_\mu J^\mu=0\), hence local conservation.
This calculation does not assume that the local transformation is a new
symmetry or that the original theory has been gauged.

\subsection{Currents are not unique}
For any local antisymmetric tensor \(U^{\nu\mu}=-U^{\mu\nu}\),
\begin{equation}
J'^\mu=J^\mu+\partial_\nu U^{\nu\mu}
\quad\Longrightarrow\quad
\partial_\mu J'^\mu=\partial_\mu J^\mu.
\label{eq:improvement}
\end{equation}
Commuting derivatives annihilates the antisymmetric term. Its integrated
charge is a surface contribution and vanishes under suitable boundary
conditions. Reversing a generator's parameter also reverses its current
and charge. These freedoms explain some harmless differences between
textbooks; they do not justify changing signs midway through a derivation.

\section{From a local current to a conserved charge}

Write \(J^\mu=(\rho,\bm j)\). In our convention,
\begin{equation}
\boxed{\partial_\mu J^\mu=\partial_t\rho+\nabla\cdot\bm j=0.}
\label{eq:continuity}
\end{equation}
There is no minus sign in this coordinate divergence. The metric enters
when raising or lowering current components: \(J_i=-J^i\).

For a fixed spatial volume \(V\), define
\begin{equation}
Q_V(t)=\int_V\dd^d x\,J^0(t,\bm x).
\end{equation}
Integrating Eq.~\eqref{eq:continuity} and using the divergence theorem,
\begin{equation}
\boxed{\frac{\dd Q_V}{\dd t}
=-\int_{\partial V}\bm j\cdot\bm n\,\dd A.}
\label{eq:chargeFlux}
\end{equation}
A decrease of charge inside the volume is exactly the outward flux.
For a spacetime slab \(B=[t_i,t_f]\times V\), the same statement is
\begin{equation}
Q_V(t_f)-Q_V(t_i)
+\int_{t_i}^{t_f}\dd t\int_{\partial V}\bm j\cdot\bm n\,\dd A=0.
\end{equation}
The final time slice has positive orientation and the initial slice has
negative orientation. The side of the slab supplies the spatial flux.

A time-independent total charge follows when the integral is finite and
the boundary flux vanishes: for example on a compact space without
boundary, with periodic boundary conditions, or on \(\R^d\) with adequate
falloff. Local conservation alone does not guarantee conservation in an
open subsystem. It also does not guarantee that the total charge of an
infinite homogeneous configuration is finite.

More geometrically, a charge on a spacelike hypersurface is
\begin{equation}
Q[\Sigma]=\int_\Sigma\dd\Sigma_\mu J^\mu.
\end{equation}
It is unchanged under deformation of the surface when there is no flux
through the remaining boundary. This is the spacetime meaning of a
conserved charge.

\section{Example 1: translations and energy-momentum}

\subsection{Translations give one current for each direction}
Suppose \(\Lag\) has no explicit spacetime dependence and the backgrounds
and boundary conditions admit translations. For
\begin{equation}
x'^\mu=x^\mu+a^\mu,\qquad \phi'^a(x')=\phi^a(x),
\end{equation}
we have \(X_{(\nu)}^\mu=\delta^\mu{}_{\nu}\),
\(\Phi_{(\nu)}^a=0\), and \(K_{(\nu)}^\mu=0\). Therefore
\begin{equation}
J_{(\nu)}^\mu=-\Th^\mu{}_{\nu},\qquad
\boxed{\partial_\mu\Th^\mu{}_{\nu}=0.}
\label{eq:translationCurrent}
\end{equation}
The \(D\) columns of the canonical energy-momentum tensor are conserved
currents. Its time row gives their densities.

\begin{remarkbox}[Fix the sign of momentum once]
With the transformation \(x'=x+a\) and the current convention
\eqref{eq:NoetherCurrent}, the direct Noether charge is
\(Q_{(\nu)}=-\int\Th^0{}_{\nu}\,\dd^d x\).
We define the conventional physical momentum by
\[
P_\nu:=\int\dd^d x\,\Th^0{}_{\nu},\qquad
P^\nu=\eta^{\nu\rho}P_\rho=(H,\bm P).
\]
Thus \(Q_{(\nu)}=-P_\nu\). Equivalently, parameterizing the same
translations by \(x'=x-a\) yields the current \(+\Th^\mu{}_{\nu}\).
The conservation law is unchanged. The charge always integrates the
\(\mu=0\) component, not a free \(\mu\) component.
\end{remarkbox}

At fixed coordinates, \(\delta\phi^a=-a^\nu\partial_\nu\phi^a\), so
\(\delta\Lag=-a^\nu\partial_\nu\Lag\). The fixed-coordinate formula
uses \(F^\mu=-a^\mu\Lag\) and gives the same answer.

\subsection{A direct divergence check}
Differentiate Eq.~\eqref{eq:Theta} and use the chain rule:
\begin{align}
\partial_\mu\Th^\mu{}_{\nu}
&=(\partial_\mu\pi_a^\mu)\partial_\nu\phi^a
 +\pi_a^\mu\partial_\mu\partial_\nu\phi^a
 -\partial_\nu\Lag\notag\\
&=-\Eul_a\partial_\nu\phi^a
 -\left(\frac{\partial\Lag}{\partial x^\nu}\right)_{\rm explicit}.
\label{eq:ThetaDivergence}
\end{align}
Hence a prescribed position-dependent coupling or source can exchange
energy-momentum with the field. Translation invariance is a property of
the complete setup, not just of the background metric.

\subsection{The real scalar: energy density and energy flow}
For
\begin{equation}
\Lag=\frac12\partial_\mu\phi\,\partial^\mu\phi-V(\phi),
\qquad \Box\phi+V'(\phi)=0,
\end{equation}
we find
\begin{equation}
\Th^\mu{}_{\nu}=\partial^\mu\phi\,\partial_\nu\phi
-\delta^\mu{}_{\nu}\Lag,
\qquad
\Th^{\mu\nu}=\partial^\mu\phi\,\partial^\nu\phi-\eta^{\mu\nu}\Lag.
\label{eq:scalarTheta}
\end{equation}
For this scalar Lagrangian \(\Th^{\mu\nu}\) is already symmetric.
The energy density and momentum density are
\begin{align}
\mathcal H=\Th^{00}
&=\frac12\dot\phi^2+\frac12|\nabla\phi|^2+V(\phi),\notag\\
\Th^{0i}&=-\dot\phi\,\partial_i\phi.
\label{eq:scalarEnergy}
\end{align}
Since the canonical momentum is \(\pi=\dot\phi\), the first formula also
follows from \(\mathcal H=\pi\dot\phi-\Lag\). Energy conservation is
\begin{equation}
\partial_t\mathcal H+\nabla\cdot(-\dot\phi\nabla\phi)=0.
\end{equation}
The energy flux is therefore \(-\dot\phi\nabla\phi\), equal to the
momentum density in this relativistic scalar example.

For a free massive field, \(V=m^2\phi^2/2\), consider
\(\phi=A\cos(\omega t-\bm k\cdot\bm x)\), with
\(\omega^2=|\bm k|^2+m^2\). Averaging over a period gives
\begin{equation}
\langle\mathcal H\rangle=\frac{A^2\omega^2}{2},\qquad
\langle\bm j_E\rangle=\frac{A^2\omega\bm k}{2},\qquad
\frac{\langle\bm j_E\rangle}{\langle\mathcal H\rangle}
=\frac{\bm k}{\omega}.
\end{equation}
The flux points in the propagation direction, and the ratio is the group
velocity. This checks both the metric sign and the physical interpretation.
The plane wave has infinite total energy on all of \(\R^d\); these are
local densities, or quantities in a periodic box.

\section{Example 2: the global phase symmetry of a complex scalar}

\subsection{An internal rotation of two real components}
Take
\begin{equation}
\Lag=\partial_\mu\phi^*\partial^\mu\phi-U(\phi^*\phi),
\qquad \phi=\frac{\phi_1+i\phi_2}{\sqrt2}.
\label{eq:complexL}
\end{equation}
The free example has \(U(z)=m^2z\). More generally, any real function of
\(z=\phi^*\phi\) preserves the constant phase rotation
\begin{equation}
\phi\longmapsto e^{i\alpha}\phi,\qquad
\phi^*\longmapsto e^{-i\alpha}\phi^*.
\end{equation}
In real components this is the rotation matrix introduced earlier:
\(\delta\phi_1=-\alpha\phi_2\), \(\delta\phi_2=\alpha\phi_1\).
The connected group is \(U(1)\simeq SO(2)\).
The displayed Lagrangian also admits the discrete reflection
\(\phi\leftrightarrow\phi^*\), so its full real-component internal
symmetry includes \(O(2)\). The continuous current comes from its
\(SO(2)\) subgroup.

The word global means that the same angle acts at every point. If
\(\alpha=\alpha(x)\), derivatives of the angle appear in the kinetic
term. The action \eqref{eq:complexL} has not yet introduced a gauge field
to compensate for them.

\subsection{Derivation and a direct conservation check}
Here \(X^\mu=0\), \(K^\mu=0\), and the field generators are
\(\Phi_\phi=i\phi\), \(\Phi_{\phi^*}=-i\phi^*\). Since
\begin{equation}
\pi_\phi^\mu=\partial^\mu\phi^*,\qquad
\pi_{\phi^*}^\mu=\partial^\mu\phi,
\end{equation}
the current is
\begin{equation}
\boxed{J^\mu=i\left[(\partial^\mu\phi^*)\phi
-\phi^*\partial^\mu\phi\right]
=\phi_1\partial^\mu\phi_2-\phi_2\partial^\mu\phi_1.}
\label{eq:U1}
\end{equation}
It is real and agrees with the phase convention \(e^{+i\alpha}\) used
on the slides. Both real components are essential.

The field equations are
\begin{equation}
\Box\phi+U'(z)\phi=0,\qquad
\Box\phi^*+U'(z)\phi^*=0.
\end{equation}
Therefore
\begin{align}
\partial_\mu J^\mu
&=i\left[(\Box\phi^*)\phi-\phi^*\Box\phi\right]\notag\\
&=i\left[-U'(z)\phi^*\phi+\phi^*U'(z)\phi\right]=0.
\end{align}
The terms containing one derivative on each field cancel before using the
field equations. The potential terms cancel because the potential is
phase invariant, not because the theory is free.

\subsection{What the charge measures}
The density and charge are
\begin{equation}
J^0=\phi_1\dot\phi_2-\phi_2\dot\phi_1,
\qquad Q=\int\dd^d x\,J^0.
\end{equation}
This is angular momentum in the \((\phi_1,\phi_2)\) plane at each
spatial point, integrated over space. Where polar variables are regular,
write \(\phi=r e^{i\theta}/\sqrt2\). Then
\begin{equation}
\Lag=\frac12(\partial r)^2+\frac12r^2(\partial\theta)^2-U(r^2/2),
\qquad J^\mu=r^2\partial^\mu\theta.
\end{equation}
The angle has no potential dependence, so its conjugate momentum is
conserved. This is the field-theory version of a cyclic angular coordinate
in mechanics.

For a free positive-frequency mode
\(\phi=Ae^{-i\omega t+i\bm k\cdot\bm x}\), our convention gives
\begin{equation}
J^0=-2\omega|A|^2.
\end{equation}
This is allowed: a charge density is not a positive probability density.
Using the inverse phase convention \(e^{-i\alpha}\) reverses the current
and charge. The energy density remains positive for the stable free
massive theory. The classical derivation by itself does not yet imply
charge quantization or a particle-number interpretation.

\section{Example 3: Lorentz transformations}

\subsection{Infinitesimal transformations and antisymmetric parameters}
Let
\begin{equation}
\Lambda^\mu{}_{\nu}=\delta^\mu{}_{\nu}+\omega^\mu{}_{\nu}
+O(\omega^2).
\end{equation}
Expanding \(\Lambda^T\eta\Lambda=\eta\) gives
\begin{equation}
\omega_{\mu\nu}:=\eta_{\mu\rho}\omega^\rho{}_{\nu}
=-\omega_{\nu\mu}.
\label{eq:omegaAntisym}
\end{equation}
It is the tensor with both indices lowered (or both raised) that is
antisymmetric. The mixed-index matrix is not generally antisymmetric
under ordinary matrix transpose; boosts give a familiar counterexample.
There are \(D(D-1)/2\) independent parameters.

The displacement is
\begin{equation}
\xi^\mu=\omega^\mu{}_{\nu}x^\nu
=\frac12\omega^{\rho\sigma}
\left(\delta^\mu{}_{\rho}x_\sigma-\delta^\mu{}_{\sigma}x_\rho\right).
\label{eq:LorentzXi}
\end{equation}
We sum over both orders of the pair \((\rho,\sigma)\), hence the factor
\(1/2\). If we sum only over independent pairs \(\rho<\sigma\), that
factor is absent. Also \(\partial_\mu\xi^\mu=0\).

\subsection{Scalar fields: orbital angular momentum}
For a scalar, \(\Delta\phi=0\), so
\(\delta\phi=-\omega^\mu{}_{\nu}x^\nu\partial_\mu\phi\).
For a Lorentz-invariant scalar Lagrangian with \(K=0\), its full current
including the parameter is
\begin{align}
\mathcal J^\mu[\omega]
&=-\Th^\mu{}_{\nu}\omega^\nu{}_{\lambda}x^\lambda\notag\\
&=\frac12\omega^{\rho\sigma}M^\mu{}_{\rho\sigma},
\qquad
\boxed{M^\mu{}_{\rho\sigma}
=x_\rho\Th^\mu{}_{\sigma}-x_\sigma\Th^\mu{}_{\rho}.}
\label{eq:orbitalM}
\end{align}
This defines the conventional normalization of the angular-momentum
current. The coefficient of \(\omega^{\rho\sigma}\) in a sum over all
ordered pairs is \(M^\mu{}_{\rho\sigma}/2\).

Assuming translation conservation, differentiation gives
\begin{equation}
\partial_\mu M^\mu{}_{\rho\sigma}
=\Th_{\rho\sigma}-\Th_{\sigma\rho}.
\label{eq:orbitalDivergence}
\end{equation}
For the standard scalar in Eq.~\eqref{eq:scalarTheta}, the tensor is
symmetric and the Lorentz current is conserved. The charges are
\begin{equation}
M_{\rho\sigma}=\int\dd^d x\,M^0{}_{\rho\sigma}.
\label{eq:LorentzCharge}
\end{equation}
Spatial moments require sufficient falloff for the factors of \(x\) as
well as for the original translation charges.

\subsection{Rotations and boosts have concrete conserved charges}
In \(3+1\) dimensions, the three rotation charges can be written
\begin{equation}
L^k=\frac12\epsilon^{kij}M_{ij},\qquad \epsilon^{123}=+1.
\end{equation}
Using raised spatial momentum components, this gives
\begin{equation}
\bm L=\int\dd^3x\,\bm x\times\bm{\mathcal P},
\qquad \mathcal P^i=\Th^{0i}.
\end{equation}
The two minus signs from lowering \(x_i\) and \(\Th^0{}_i\) cancel.

For a symmetric conserved stress tensor, define the three boost charges
with raised indices,
\begin{equation}
B^i:=M^{0i}=tP^i-\int\dd^3x\,x^i\mathcal H.
\label{eq:boost}
\end{equation}
Energy conservation and an integration by parts give
\begin{equation}
\frac{\dd}{\dd t}\int x^i\mathcal H\,\dd^3x
=\int\Th^{i0}\,\dd^3x=P^i,
\end{equation}
so \(\dd B^i/\dd t=0\). If \(H\ne0\), the center of energy
\(X_E^i=H^{-1}\int x^i\mathcal H\,\dd^3x\) consequently satisfies
\(\dot X_E^i=P^i/H\). Conservation of boosts expresses uniform motion of
the center of energy. A conserved charge may contain explicit time
dependence that cancels the evolution of the fields.

\subsection{Vectors and other fields also rotate their components}
A covector transforms by
\begin{equation}
A'_\mu(x')=(\Lambda^{-1})^\nu{}_{\mu}A_\nu(x),
\end{equation}
so the one-form \(A'_\mu(x')\dd x'^\mu=A_\mu(x)\dd x^\mu\) is preserved.
Its two variations are
\begin{align}
\Delta A_\mu&=-\omega^\nu{}_{\mu}A_\nu,\notag\\
\delta A_\mu&=-\xi^\nu\partial_\nu A_\mu
             -A_\nu\partial_\mu\xi^\nu.
\label{eq:covectorVariation}
\end{align}
The second expression is minus the Lie derivative of a one-form. Its
first term moves the argument; its second changes the component basis.
For a contravariant vector the corresponding component term has the
opposite index action:
\(\delta A^\mu=-\xi^\nu\partial_\nu A^\mu
+A^\nu\partial_\nu\xi^\mu\).

For a general multiplet let
\begin{equation}
\phi'^a(\Lambda x)=S(\Lambda)^a{}_b\phi^b(x),\qquad
S(\Lambda)=1+\frac12\omega^{\rho\sigma}\mathsf S_{\rho\sigma}
+O(\omega^2),
\label{eq:LorentzRepresentation}
\end{equation}
where \(\mathsf S_{\rho\sigma}=-\mathsf S_{\sigma\rho}\) are matrices
acting on field components. This definition absorbs any factors of \(i\)
into the matrices; conventions with explicit \(i\) must change the
matrix definitions accordingly. Spinors use a representation of the
Lorentz group's double cover. These finite-dimensional component
representations need not be unitary.

The component variation contributes a spin current
\begin{equation}
s^\mu{}_{\rho\sigma}
=\pi_a^\mu(\mathsf S_{\rho\sigma})^a{}_b\phi^b.
\end{equation}
For a strictly invariant action, the full Lorentz current is
\begin{equation}
\boxed{M^\mu{}_{\rho\sigma}
=x_\rho\Th^\mu{}_{\sigma}-x_\sigma\Th^\mu{}_{\rho}
+s^\mu{}_{\rho\sigma}.}
\label{eq:totalM}
\end{equation}
Noether's theorem now implies, on shell,
\begin{equation}
\partial_\mu s^\mu{}_{\rho\sigma}
=\Th_{\sigma\rho}-\Th_{\rho\sigma}.
\label{eq:spinIdentity}
\end{equation}
A nonsymmetric canonical tensor is therefore compatible with Lorentz
invariance: the spin current supplies the missing contribution.

\begin{warningbox}[The symmetry condition on the stress tensor needs a qualifier]
For the purely orbital current, after translation conservation is imposed,
Eq.~\eqref{eq:orbitalDivergence} vanishes exactly when the stress tensor is
symmetric. For fields with a spin contribution, it is the \emph{total}
current \eqref{eq:totalM} that is conserved. The canonical stress tensor
need not be symmetric.
\end{warningbox}

In ordinary Poincaré-invariant theories, a suitable improvement can absorb
the spin contribution into a symmetric stress tensor. This is the
Belinfante construction. The improvement changes local densities by
superpotential terms and leaves the integrated Poincare charges unchanged
when the corresponding boundary terms vanish. Maxwell theory from Lecture
6 is an example where the distinction between a canonical tensor and a
symmetric, gauge-invariant tensor matters. Its construction can be developed
after the present scalar examples.

\section{How to apply the theorem}
For each proposed symmetry:
\begin{enumerate}
\item Specify the action, field components, background data, and boundary
conditions. State the metric and generator conventions.
\item Write \(\xi^\mu=\varepsilon^I X_I^\mu\) and
\(\Delta\phi^a=\varepsilon^I\Phi_I^a\); calculate
\(\chi_I^a=\Phi_I^a-X_I^\nu\partial_\nu\phi^a\).
\item Check the local symmetry identity off shell and determine \(K_I^\mu\).
Equivalently, compute the fixed-coordinate boundary term \(F_I^\mu\).
\item Compute \(\pi_a^\mu\) and then
\(J_I^\mu=\pi_a^\mu\Phi_I^a-\Th^\mu{}_{\nu}X_I^\nu-K_I^\mu\).
\item Check its divergence with the field equations. Integrate \(J_I^0\)
and state the boundary-flux condition needed to obtain a conserved charge.
\end{enumerate}

\section{Questions to think about}
\begin{enumerate}
\item Translate a localized Gaussian profile by \(+a\). Compute
\(\delta\phi\) and \(\Delta\phi\), and explain why they differ.
\item Derive Eqs.~\eqref{eq:derivativeVariation} and \eqref{eq:Jacobian}.
Why can one commute \(\delta\), but not \(\Delta\), with a derivative?
\item For the massless scalar with \(\Lag=(\partial\phi)^2/2\), find the
current for \(\phi\mapsto\phi+\varepsilon\). What term in a potential
would break this symmetry?
\item Let \(\Lag=(\partial\phi)^2/2-m^2(x)\phi^2/2\). Use
Eq.~\eqref{eq:ThetaDivergence} to determine the local rate of
energy-momentum exchange with the prescribed background.
\item Compute the scalar energy flux for a wave moving in the negative
\(x\)-direction. Check the sign of \(\Th^{01}\).
\item Add \(-\lambda(\phi^*\phi)^2\) to the complex scalar Lagrangian.
Show that the phase current is unchanged and remains conserved.
\item Instead add \(-\kappa(\phi^2+\phi^{*2})/2\), for real \(\kappa\).
Find \(\partial_\mu J^\mu\) on shell and identify the surviving discrete
phase symmetry.
\item Starting from \(J^0=\phi_1\dot\phi_2-\phi_2\dot\phi_1\), show that
\(\{\phi_1,Q\}=-\phi_2\) and \(\{\phi_2,Q\}=\phi_1\), using equal-time
canonical Poisson brackets. In what sense does the charge generate the
symmetry?
\item Prove the boost-charge conservation law \eqref{eq:boost}. Which
surface term and which symmetry of \(\Th\) are needed?
\item Derive the covector variation \eqref{eq:covectorVariation} from
invariance of \(A_\mu\dd x^\mu\). Identify its component and orbital parts.
\item Explain why an antisymmetric improvement preserves local
conservation but may change a charge in a region with a boundary.
\item In mechanics, take \(L=m\dot q^2/2\) and the boost
\(\delta q=\varepsilon t\) at fixed time. Show that
\(\delta L=\varepsilon\,\dd(mq)/\dd t\), and use the boundary term to
find the conserved quantity \(pt-mq\).
\end{enumerate}

\section*{Summary}
\addcontentsline{toc}{section}{Summary}
The geometric field variation and the fixed-coordinate variation are
related by \(\delta\phi^a=\Delta\phi^a-\xi^\mu\partial_\mu\phi^a\).
This distinction, together with the Jacobian of the coordinate map, gives
the action's boundary term. For a continuous global variational symmetry,
\begin{equation*}
J_I^\mu=\pi_a^\mu\Phi_I^a-\Th^\mu{}_{\nu}X_I^\nu-K_I^\mu,
\qquad \partial_\mu J_I^\mu=-\Eul_a\chi_I^a.
\end{equation*}
The right-hand side vanishes on classical solutions. A finite charge is
conserved when the associated boundary flux vanishes.

Translations give energy-momentum conservation; a complex scalar's phase
rotation gives angular momentum in its internal two-component plane;
Lorentz transformations give rotations and boosts. Scalars have a purely
orbital Lorentz current, while fields transforming in nontrivial component
representations also have a spin current. The current's sign and
normalization must track the chosen symmetry parameter throughout.

\smallskip
\noindent\textit{Further reading:} David Tong,
\href{https://www.damtp.cam.ac.uk/user/tong/qft.htm}{Quantum Field Theory},
Section 1.3, for another presentation of Noether currents, translations,
and Lorentz transformations. The derivation here follows the supplied
Lecture 7 slides, with its conventions stated explicitly.
\end{document}
